% GENERATED FILE. Edit canonical lessons, then run npm run generate:books.
% canonical-source-hash: fnv1a64:6a11649fcbb90611
% canonical-lessons: TE-C156-kalapattika, TE-C156-sucana, TE-C156-adde, TE-C156-prakatana, TE-C156-rasidu, TE-R156-first-pass-saying-no-and-looking-back, TE-R156-second-pass-looking-ahead-and-reading-signs

\chapter{A Timetable, a Notice, the Rent, an Advert, a Receipt}
\label{ch:te-a-timetable-a-notice-the-rent-an-advert-a-receipt}
\hlchaptermodality{156}

\begin{chapteropening}
\textbf{By the end of this chapter:} \emph{I can read a timetable, a notice, an advert and a receipt, and act on what they say.}

Complete the last lesson of chapter 156: I can read a timetable, a notice, an advert and a receipt, and act on what they say.
\end{chapteropening}

\section[\texorpdfstring{\te{కాలపట్టిక} (kālapaṭṭika)}{కాలపట్టిక (kālapaṭṭika)}]{\te{కాలపట్టిక} (kālapaṭṭika) — a timetable}

\label{lesson:TE-C156-kalapattika}



\begin{quote}
Before the new one: say the Telugu for to hope, then the Telugu for to look forward to.
\end{quote}

\subsection*{You'll want to know: \te{కాలపట్టిక}}
\textbf{\te{కాలపట్టిక}} — \emph{kālapaṭṭika} — \textquotedblleft{}a timetable\textquotedblright{}. Read the board: \textbf{\te{బస్సు}: \te{ఉదయం} \te{ఎనిమిది} \te{గంటలు}} — \emph{bassu: udayaṁ enimidi gaṇṭalu} — \textquotedblleft{}Bus: 8 a.m.\textquotedblright{} The bus leaves at eight, so be there before eight.

\begin{grammarlens}[title={What you've built}]
One more word for this chapter.
\end{grammarlens}

\subsection*{Your turn}
\begin{itemize}
\raggedright
  \item \emph{Say it:} \emph{kālapaṭṭika}
  \item \emph{Say it:} \emph{kālapaṭṭika}, once more
  \item \emph{Say it:} say \emph{kālapaṭṭika}
  \item \emph{Recall:} say the Telugu for to hope, then the Telugu for to look forward to, then say \emph{kālapaṭṭika} again
\end{itemize}

\begin{tcolorbox}[breakable,colback=teal!4,colframe=teal!35!black,title={Before you move on}]
What does \te{కాలపట్టిక} mean? (\textquotedblleft{}A timetable\textquotedblright{}.) Say it once more.
\end{tcolorbox}
\section[\texorpdfstring{\te{సూచన} (sūcana)}{సూచన (sūcana)}]{\te{సూచన} (sūcana) — an instruction, a note (on a sign)}

\label{lesson:TE-C156-sucana}



\begin{quote}
Before the new one: say the Telugu for to look forward to, then the Telugu for a timetable.
\end{quote}

\subsection*{You'll want to know: \te{సూచన}}
\textbf{\te{సూచన}} — \emph{sūcana} — \textquotedblleft{}an instruction, a note (on a sign)\textquotedblright{}. Read the notice: \textbf{\te{ఇవాళ} \te{దుకాణం} \te{మూసి} \te{ఉంది}} — \emph{ivāḷa dukāṇaṁ mūsi undi} — \textquotedblleft{}The shop is closed today.\textquotedblright{} Come back on another day.

\begin{grammarlens}[title={What you've built}]
One more word for this chapter.
\end{grammarlens}

\subsection*{Your turn}
\begin{itemize}
\raggedright
  \item \emph{Say it:} \emph{sūcana}
  \item \emph{Say it:} \emph{sūcana}, once more
  \item \emph{Say it:} say \emph{sūcana}
  \item \emph{Recall:} say the Telugu for to look forward to, then the Telugu for a timetable, then say \emph{sūcana} again
\end{itemize}

\begin{tcolorbox}[breakable,colback=teal!4,colframe=teal!35!black,title={Before you move on}]
What does \te{సూచన} mean? (\textquotedblleft{}An instruction, a note (on a sign)\textquotedblright{}.) Say it once more.
\end{tcolorbox}
\section[\texorpdfstring{\te{అద్దె} (adde)}{అద్దె (adde)}]{\te{అద్దె} (adde) — rent}

\label{lesson:TE-C156-adde}



\begin{quote}
Before the new one: say the Telugu for a timetable, then the Telugu for an instruction.
\end{quote}

\subsection*{You'll want to know: \te{అద్దె}}
\textbf{\te{అద్దె}} — \emph{adde} — \textquotedblleft{}rent\textquotedblright{}. The money you pay to live in a house: \textbf{\te{ఇంటి} \te{అద్దె}} — \emph{iṇṭi adde} — \textquotedblleft{}the rent for the house\textquotedblright{}.

\begin{grammarlens}[title={What you've built}]
One more word for this chapter.
\end{grammarlens}

\subsection*{Your turn}
\begin{itemize}
\raggedright
  \item \emph{Say it:} \emph{adde}
  \item \emph{Say it:} \emph{adde}, once more
  \item \emph{Say it:} say \emph{adde}
  \item \emph{Recall:} say the Telugu for a timetable, then the Telugu for an instruction, then say \emph{adde} again
\end{itemize}

\begin{tcolorbox}[breakable,colback=teal!4,colframe=teal!35!black,title={Before you move on}]
What does \te{అద్దె} mean? (\textquotedblleft{}Rent\textquotedblright{}.) Say it once more.
\end{tcolorbox}
\section[\texorpdfstring{\te{ప్రకటన} (prakaṭana)}{ప్రకటన (prakaṭana)}]{\te{ప్రకటన} (prakaṭana) — an advertisement, an announcement}

\label{lesson:TE-C156-prakatana}



\begin{quote}
Before the new one: say the Telugu for an instruction, then the Telugu for rent.
\end{quote}

\subsection*{You'll want to know: \te{ప్రకటన}}
\textbf{\te{ప్రకటన}} — \emph{prakaṭana} — \textquotedblleft{}an advertisement, an announcement\textquotedblright{}. Read the advert: \textbf{\te{ఇల్లు} \te{అద్దెకు} \te{ఉంది}} — \emph{illu addeku undi} — \textquotedblleft{}House for rent.\textquotedblright{}

\begin{grammarlens}[title={What you've built}]
One more word for this chapter.
\end{grammarlens}

\subsection*{Your turn}
\begin{itemize}
\raggedright
  \item \emph{Say it:} \emph{prakaṭana}
  \item \emph{Say it:} \emph{prakaṭana}, once more
  \item \emph{Say it:} say \emph{prakaṭana}
  \item \emph{Recall:} say the Telugu for an instruction, then the Telugu for rent, then say \emph{prakaṭana} again
\end{itemize}

\begin{tcolorbox}[breakable,colback=teal!4,colframe=teal!35!black,title={Before you move on}]
What does \te{ప్రకటన} mean? (\textquotedblleft{}An advertisement, an announcement\textquotedblright{}.) Say it once more.
\end{tcolorbox}
\section[\texorpdfstring{\te{రసీదు} (rasīdu)}{రసీదు (rasīdu)}]{\te{రసీదు} (rasīdu) — a receipt}

\label{lesson:TE-C156-rasidu}



\begin{quote}
Before the new one: say the Telugu for rent, then the Telugu for an advertisement.
\end{quote}

\subsection*{You'll want to know: \te{రసీదు}}
\textbf{\te{రసీదు}} — \emph{rasīdu} — \textquotedblleft{}a receipt\textquotedblright{}. Ask for it when you pay: \textbf{\te{రసీదు} \te{ఇస్తారా}?} — \emph{rasīdu istārā?} — \textquotedblleft{}Could you give me the receipt?\textquotedblright{} Read two words again: \textbf{\te{ఛాయ}} (\emph{chāya}, shade) is written with \te{ఛ}, and \textbf{\te{ఝరి}} (\emph{jhari}, a stream) with \te{ఝ}.

\begin{grammarlens}[title={What you've built}]
One more word for this chapter.
\end{grammarlens}

\subsection*{Your turn}
\begin{itemize}
\raggedright
  \item \emph{Say it:} \emph{rasīdu}
  \item \emph{Say it:} \emph{rasīdu}, once more
  \item \emph{Say it:} say \emph{rasīdu}
  \item \emph{Recall:} say the Telugu for rent, then the Telugu for an advertisement, then say \emph{rasīdu} again
\end{itemize}

\begin{tcolorbox}[breakable,colback=teal!4,colframe=teal!35!black,title={Before you move on}]
What does \te{రసీదు} mean? (\textquotedblleft{}A receipt\textquotedblright{}.) Say it once more.
\end{tcolorbox}
\section[(dialogue)]{Saying no and looking back — a first pass}

\label{lesson:TE-R156-first-pass-saying-no-and-looking-back}



\begin{quote}
Say the words for an advertisement and a receipt.
\end{quote}

\subsection*{Your turn}
Say these aloud:

\begin{itemize}
\raggedright
  \item nobody, never, nowhere, anyone, anything
  \item recently, already, last, once, to spend
\end{itemize}

\begin{tcolorbox}[breakable,colback=teal!4,colframe=teal!35!black,title={Before you move on}]
Nobody? (\textbf{\te{ఎవరూ}}, \emph{evarū}.) Never? (\textbf{\te{ఎప్పుడూ}}, \emph{eppuḍū}.) Nowhere? (\textbf{\te{ఎక్కడా}}, \emph{ekkaḍā}.) Anyone? (\textbf{\te{ఎవరైనా}}, \emph{evarainā}.) Anything? (\textbf{\te{ఏమైనా}}, \emph{ēmainā}.) Recently? (\textbf{\te{ఇటీవల}}, \emph{iṭīvala}.) Already? (\textbf{\te{ఇప్పటికే}}, \emph{ippaṭikē}.) Last? (\textbf{\te{గత}}, \emph{gata}.) Once? (\textbf{\te{ఒకసారి}}, \emph{okasāri}.) To spend? (\textbf{\te{గడుపు}}, \emph{gaḍupu}.)
\end{tcolorbox}
\section[(dialogue)]{Looking ahead and reading signs — a second pass}

\label{lesson:TE-R156-second-pass-looking-ahead-and-reading-signs}



\begin{quote}
Say the words for next and from now on.
\end{quote}

\subsection*{Your turn}
Say these aloud:

\begin{itemize}
\raggedright
  \item next, from now on, the future, to hope, to look forward to
  \item a timetable, an instruction, rent, an advertisement, a receipt
\end{itemize}

\begin{tcolorbox}[breakable,colback=teal!4,colframe=teal!35!black,title={Before you move on}]
Next? (\textbf{\te{వచ్చే}}, \emph{vaccē}.) From now on? (\textbf{\te{ఇకపై}}, \emph{ikapai}.) The future? (\textbf{\te{భవిష్యత్తు}}, \emph{bhaviṣyattu}.) To hope? (\textbf{\te{ఆశించు}}, \emph{āśiñcu}.) To look forward to? (\textbf{\te{ఎదురుచూడు}}, \emph{edurucūḍu}.) A timetable? (\textbf{\te{కాలపట్టిక}}, \emph{kālapaṭṭika}.) An instruction? (\textbf{\te{సూచన}}, \emph{sūcana}.) Rent? (\textbf{\te{అద్దె}}, \emph{adde}.) An advertisement? (\textbf{\te{ప్రకటన}}, \emph{prakaṭana}.) A receipt? (\textbf{\te{రసీదు}}, \emph{rasīdu}.)
\end{tcolorbox}
